\section{Definiciones}
\label{sec:definiciones}

Esta sección es normativa. Un término definido aquí tiene ese significado y
ningún otro en el resto del documento, y cada definición nombra el sitio del
código o del esquema donde el objeto existe de verdad. Las definiciones están
numeradas para poder citarlas.

\subsection{Nota sobre el vocabulario}

Los nombres técnicos se conservan en inglés. No es un descuido de redacción: son
los identificadores del código y de las columnas, y el documento se usa para
discutir sobre ellos. Donde el castellano aporta claridad se usa la traducción
como glosa, nunca como sustituto.

\subsection{Los objetos del registro}

\term{Protein}{Una secuencia de aminoácidos con un accession de UniProt. Es la
entidad sobre la que se predice y la unidad sobre la que se estratifica.}

\term{GO term}{Un nodo de la Gene Ontology, identificado por su \cod{go\_id}
textual. Pertenece a exactamente uno de los tres aspectos (D\ref{def:aspect}).
Un término existe siempre dentro de un \emph{ontology snapshot}: el mismo
\cod{go\_id} es una fila distinta en dos snapshots, y confundir las dos es la
forma habitual de comparar dos grafos distintos creyendo que es uno.}

\term{Annotation}{La afirmación de que una proteína tiene una función, con un
código de evidencia que dice cómo se supo. Vive en \cod{protein\_go\_annotation}
y pertenece a un \emph{annotation set}.}

\term{Annotation set}{El contenido de una release concreta de GOA, cargada
completa. Es el objeto que la campaña fecha: una release nombra un fichero y no
fecha nada por sí misma, de modo que la fecha de publicación se almacena aparte.}

\term{Evidence code}{La etiqueta de tres letras que clasifica el soporte de una
anotación. La campaña distingue dos grupos que importan para todo lo demás: los
códigos \emph{experimentales} y de \emph{alto rendimiento} (\cod{EXP}, \cod{IDA},
\cod{IPI}, \cod{IMP}, \cod{IGI}, \cod{IEP}, \cod{HTP}, \cod{HDA}, \cod{HMP},
\cod{HGI}, \cod{HEP}, más \cod{IC} y \cod{TAS}), y el resto, que es
mayoritariamente inferencia computacional (\cod{IEA}). La distinción es la que
usan CAFA y LAFA para construir su verdad de referencia, y aquí se usa en dos
sitios distintos que conviene no confundir: decide qué anotaciones pueden
\emph{donar} (D\ref{def:bank}) y decide qué anotaciones cuentan como verdad al
\emph{evaluar}.}

\term{Query}{Una proteína para la que se pide predicción en una corrida. El
conjunto de queries de la campaña es un \cod{query\_set} fijo.}

\term{Donor}{Una proteína del banco cuya anotación se transfiere a una query
porque su vector está cerca. Es el mecanismo entero del método base: no se
predice una función, se copia la de un vecino.}

\term{Candidate}{Un par (query, GO term) propuesto por el sistema, con una
puntuación. La columna \cod{go\_prediction.ref\_protein\_accession} es
\cod{NOT NULL}, de modo que el esquema prohíbe un candidato sin donante en lugar
de limitarse a registrar que nadie escribió uno.}

\term{Arm}{Una configuración completa, ejecutada de principio a fin, que produce
un conjunto de predicciones evaluable. Es la unidad que se compara con otra. Dos
arms difieren en al menos un parámetro, y nombrar cuál es precisamente el
problema que resuelve la noción de \emph{nivel} (D\ref{def:nivel}).}

\term{Prediction set}{El producto de un arm: el conjunto de candidatos con su
procedencia. Lleva en \cod{meta} los parámetros con los que se generó.}

\term{Evaluation set}{Un par de annotation sets que define una ventana temporal:
el de la cota inferior, del que pueden venir los donantes, y el de la cota
superior, que aporta la verdad de referencia. Nada más es una ventana.}

\term{Evaluation result}{Una fila con las métricas de un prediction set contra un
evaluation set. Es el único sitio donde vive un número publicable.}

\subsection{El marco}

\term[def:marco]{Frame}{El conjunto de condiciones bajo las que un número se
lee, y que por tanto deben coincidir para que dos números sean comparables. Hoy
son seis campos: el evaluation set, el pivot snapshot, el conjunto de
information accretion, la etiqueta de ventana temporal, el número máximo de
términos y la distancia máxima. La sección~\ref{sec:marco} argumenta que dos de
ellos están mal elegidos y que faltan dos.}

\term{Window}{El intervalo temporal de un evaluation set, escrito como
\emph{inferior} $\rightarrow$ \emph{superior}. La campaña selecciona sobre
$220 \rightarrow 227$. La cota inferior es la única release de la que un donante
puede venir sin leer el futuro, y de ahí que el frame fije el banco.}

\term{Pivot snapshot}{El ontology snapshot bajo el que se resuelven los términos
al comparar predicción y verdad. Importa porque la ontología cambia entre
releases: un término obsoleto o reubicado cambia de posición en el grafo, y con
ella cambia la propagación a ancestros.}

\term{Information accretion}{El peso por término que usa CAFA para que acertar
un término raro valga más que acertar uno trivialmente frecuente. Se calcula
sobre el corpus de la cota inferior y se congela en un
\cod{information\_accretion\_set}. Es el \emph{w} de \cod{fmax\_w}.}

\term{Frame seal, o frame digest}{El resumen criptográfico de los campos del
frame, almacenado en cada evaluation result. Es el único sello real: dos
resultados son comparables si y solo si su digest coincide. La columna llamada
\cod{frame}, que contiene el texto \cod{internal}, es una etiqueta del arnés de
ejecución y no puede sostener esa garantía.}

\term{Leak, o fuga}{Que un número use, directa o indirectamente, información
posterior a la cota superior de su ventana. La campaña la impide de tres
maneras, en orden de fuerza: por ausencia, no cargando en la base las releases
que no debe mirar; por guardia, negándose a evaluar una codificación cuyo corte
de entrenamiento cae dentro de la ventana; y por protocolo, que es la más débil
y la única que depende de la disciplina de quien opera.}

\subsection{Las diez decisiones}
\label{sec:diez}

El sistema se describe como una cadena de diez decisiones ordenadas. Cada una
es un \emph{nodo} (D\ref{def:nodo}) y tiene una pregunta propia. El orden es el
del flujo de datos: una decisión de etapa $n$ solo tiene sentido si las de etapa
menor ya están tomadas.

\term[def:substrate]{Substrate (etapa 1). El espacio vectorial en el que se
define la palabra \emph{vecino}}{Es la configuración de codificación completa, no
solo el nombre del modelo. Once campos la determinan: \cod{model\_name},
\cod{model\_backend}, \cod{layer\_indices}, \cod{layer\_agg}, \cod{pooling},
\cod{normalize}, \cod{normalize\_residues}, \cod{max\_length},
\cod{use\_chunking}, \cod{chunk\_size} y \cod{chunk\_overlap}. Cambiar cualquiera
cambia las distancias, y con ellas cambia quién es vecino de quién, qué donantes
se recuperan y qué se predice. Por eso el substrate es la primera decisión
sustantiva y no un detalle de implementación.

Es útil decir qué \emph{no} es. No es un clasificador: nada se entrena aquí. Es
solo la función que convierte una secuencia en un vector. Todo lo que el método
base hace después consiste en medir distancias en ese espacio y copiar
anotaciones de los más cercanos.

Nada en el frame nombra una representación, de modo que el substrate nunca está
fijado por la ventana: si solo hay un nivel instanciado, es un valor que nadie
eligió y no uno que el marco impuso.}

\term[def:bank]{Bank (etapa 2). Quien tiene derecho a donar, y qué se le permite
decir}{Es la decisión que el resto del documento llama \emph{banco}. Se compone
de dos mitades que no pueden decidirse por separado, y por eso forman un solo
nodo:

\begin{itemize}[leftmargin=1.4em,itemsep=0.2em]
  \item El \textbf{corpus}: de qué annotation set salen los donantes. Lo fija el
        frame, porque la cota inferior de la ventana es la única release de la
        que se puede donar sin leer el futuro.
  \item La \textbf{política de donantes}: qué subconjunto de ese corpus puede
        donar de hecho. Son \cod{donor\_reviewed\_only}, la lista
        \cod{donor\_evidence\_codes} y la lista \cod{donor\_exclusions} de
        prefijos de referencia excluidos.
\end{itemize}

\begin{aviso}{Tres interruptores que la prosa atribuye al banco y el código al retriever}
\cod{exclude\_self\_neighbour}, \cod{expand\_votes\_to\_ancestors} y
\cod{aspect\_separated\_knn} modifican el acto de donar y es natural leerlos como
política de banco. En el árbol no lo son: \cod{\_BANK\_FIELDS} contiene solo
\cod{bank\_source}, \cod{bank\_version}, \cod{donor\_reviewed\_only},
\cod{donor\_evidence\_codes} y \cod{donor\_exclusions}, mientras los tres booleanos
están en \cod{\_RETRIEVER\_FIELDS} bajo un comentario que dice literalmente
\emph{quién puede ser vecino, que es este nodo y no el banco}.

\textbf{La consecuencia hay que decirla entera.} El eje que la campaña declaró
cerrado como \emph{banco y donante} se midió sobre esos tres booleanos más la
política de donantes. Los tres booleanos cierran, en el código, el nodo
\cod{retriever}; el nodo \cod{bank} sigue en \cod{unpowered} porque su único
nivel puntuado es el corpus 220. Un lector no puede saber hoy si lo cerrado fue
el banco o el retriever, y el nodo \cod{bank} solo sale de \cod{unpowered}
cargando un segundo corpus, que es exactamente la escalera de recencia. Reconciliar
la agrupación en prosa con la agrupación implementada es una decisión pendiente,
no una errata.
\end{aviso}

Son un nodo porque un corpus y el filtro que se le aplica definen un banco entre
los dos, y porque un grupo de campos inseparables se sostiene con la firmeza de
su miembro más flojo: si el corpus está fijado por el marco pero la política está
vacía, el nodo entero es un valor heredado.

La distinción con el substrate, que es la confusión más fácil de cometer, se dice
en una frase: \textbf{el substrate decide quién está cerca; el bank decide a
quién de los que están cerca se le hace caso}. Un banco más permisivo admite
donantes cuya anotación se apoya en inferencia computacional, lo que aumenta la
cobertura y baja la fiabilidad por donante. Un banco restringido a evidencia
experimental hace lo contrario. Cuál de los dos gana no es evidente a priori y
es precisamente lo que la campaña mide.

Una propiedad del banco que conviene tener presente porque contamina varias
lecturas: bajo el protocolo temporal, la query se recupera a sí misma desde el
banco de la release anterior si ya estaba anotada entonces. Ese \emph{self hit}
está a distancia cero y es el donante más cercano de casi todas las proteínas
con anotación previa. \cod{exclude\_self\_neighbour} es el interruptor que lo
retira, y varias definiciones de la sección~\ref{sec:unidad} se resuelven
explícitamente sobre donantes de secuencia distinta a la query por esta razón.}

\term[def:retriever]{Retriever (etapa 3). Cómo se extraen candidatos del banco, y
a qué profundidad}{Dos parámetros y un cuidado. La \textbf{métrica} de distancia
(\cod{cosine}) y la \textbf{profundidad} de recuperación
(\cod{limit\_per\_entry}), que es cuántos vecinos se traen por query. El cuidado
es que la profundidad son en realidad dos cantidades distintas que el sistema no
debe mezclar: el número de vecinos que se \emph{recuperan}, que cambia qué
candidatos existen y exige volver a ejecutar la recuperación, y el
\emph{corte de evaluación} (\cod{max\_k\_position}), que trunca una lista ya
recuperada y no exige volver a calcular nada. Un guardia específico se niega a
comparar dos resultados cuyas profundidades sean cantidades distintas.}

\term{Generator (etapa 4). Si algún candidato llega sin donante}{Mide si existe
una vía de generación de candidatos que no sea la transferencia por vecindad,
típicamente anotación por dominios de InterPro. Hoy la respuesta es que no
existe, y el esquema lo garantiza en lugar de constatarlo.}

\term{Scoring (etapa 5). Qué ponderación convierte un candidato en una
puntuación}{La fórmula que asigna el número. En el estado actual todos los
números del registro los produjo la vía de respaldo,
$1 - \text{distancia coseno}/2$, que no está registrada como configuración. El
nodo existe para que una ponderación alternativa sea comparable con esa.}

\term{Features (etapa 6). Qué rasgos por candidato entran en un modelo}{El
conjunto de columnas que describe un candidato: identidad de secuencia con el
donante, relación taxonómica, information accretion del término, estadísticos de
la votación de vecinos, y demás. El esquema de rasgos está versionado y agrupa
21 familias sobre 20 conjuntos de columnas independientes. La ablación correcta
retira \emph{columnas}, nunca familias, porque una de las familias (\cod{knn}) es
la unión disjunta de otras dos (\cod{knn\_distance} y \cod{knn\_vote}); las otras
dieciocho son disjuntas dos a dos, y de ahí que 21 familias den 20 conjuntos.}

\term{Reranking (etapa 7). Si un modelo reordena los candidatos}{El paso que
convierte el sistema de un método de recuperación en un método aprendido. Toma
los rasgos de la etapa 6 y reordena. Es el objeto de la segunda pregunta de
investigación del proyecto y hoy no tiene ninguna instancia.}

\term{Combination (etapa 8). Cómo se fusionan dos o más flujos en una respuesta}{
Relevante solo cuando existe más de una vía de generación, es decir, cuando el
nodo generator deja de estar bloqueado.}

\term[def:routing]{Routing (etapa 9). Qué flujo responde a qué panel}{La decisión
de encaminar configuraciones distintas a regiones distintas de la población, en
lugar de elegir una sola configuración para todo. Es lo que el proyecto llama
\emph{canalización por estratificación masiva}, y la sección~\ref{sec:estadistico}
explica por qué es la decisión con más riesgo estadístico de las diez.}

\subsection{Nodo, nivel, arista y firmeza}

\term[def:nodo]{Node}{Una decisión sobre un campo, o sobre un grupo de campos que
no pueden decidirse por separado. Las diez de la sección~\ref{sec:diez} son todas
las que el sistema reconoce, y la lista es fija en el código: una prueba recorre
la especificación, resuelve el constructor de cada nodo por su nombre, comprueba
que hay diez y comprueba que a los diez se les entrega el umbral declarado.}

\term[def:nivel]{Level}{\textbf{Un valor junto con todos los campos que varían
entre las filas con las que se compara.} Esta es la definición más importante del
documento y la fuente del defecto que más veces ha reaparecido en el proyecto.

El motivo es que un nombre corto es casi siempre insuficiente. Si se comparan dos
arms que difieren en el modelo y \emph{también} en la profundidad de
recuperación, entonces \cod{ankh\_large} no nombra un nivel: nombra dos arms
distintos que comparten el nombre del modelo. La comparación parece de un solo
campo y no lo es.

De aquí se sigue una propiedad incómoda pero real: \textbf{la aridad de un nivel
se decide en el momento de la lectura, no en el momento de la declaración}. Sobre
el registro completo hacen falta cinco campos para nombrar un nivel; restringido
a un solo frame seal, \cod{donor\_policy} pasa a ser constante y bastan cuatro.
La consecuencia práctica es que \emph{ninguna cadena de texto almacenada puede
nombrar un nivel}, y por tanto un umbral debe declararse como una tupla de campos
más valores, nunca como un nombre ya compuesto.}

\term{Edge}{Lo que el sistema sabe sobre un nodo: si el artefacto que hace falta
existe, cuántos niveles se han instanciado, cuántos están disponibles, cuántos se
han puntuado, cuántos resultados hay, si el marco fuerza el valor, cuál es el
umbral declarado y si la comparación lo superó.}

\term[def:firmeza]{Strength}{La palabra que dice con qué firmeza se sostiene una
decisión. Hay cinco y se prueban en un orden fijo, que es el argumento del
diseño:

\begin{center}
\begin{tabularx}{\linewidth}{@{}c L{3.2cm} X@{}}
\toprule
\# & palabra & se devuelve cuando\\
\midrule
1 & \cod{blocked}    & el artefacto no existe, o no hay ningún nivel instanciado\\
2 & \cod{chosen} o \cod{inherited} & hay un solo nivel. \cod{chosen} si el marco lo fuerza, \cod{inherited} si nadie lo eligió\\
3 & \cod{unpowered}  & hay dos niveles o más, pero menos de dos puntuados\\
4 & \cod{chosen}     & hay poder, pero no hay umbral declarado, o el umbral no pudo responder\\
5 & \cod{measured}   & hay poder, hay umbral, y la comparación lo superó. Si no lo supera, \cod{chosen}\\
\bottomrule
\end{tabularx}
\end{center}

Dos propiedades del orden importan aguas abajo. \textbf{El poder se prueba antes
que la evidencia}: ninguna diferencia, por grande que sea, puede convencer a un
nodo de que hubo una medida si la comparación no tenía forma de resolver nada.
Y \textbf{\cod{chosen} es un sumidero}: cuatro situaciones distintas caen en esa
palabra, de modo que un lector que solo vea la palabra no puede distinguir una
decisión sin umbral de una que tenía umbral y no lo superó. Publicar el umbral y
el resultado de la comparación junto al nodo es un cambio pequeño y resuelve eso.}

\term[def:suelo]{Floor}{El nivel incumbente contra el que se compara, declarado
\emph{antes} de medir. Es lo único que hace falta desde fuera de las tablas de
resultados para que un nodo pase de \cod{chosen} a \cod{measured}, y por eso es
el punto donde una campaña se vuelve falsable o deja de serlo.}

\term[def:haz]{Survivor set, o haz}{El conjunto de niveles que ninguna celda con
poder llegó a batir. Es el resultado correcto cuando un eje no sella: no un
ganador, ni un empate, sino los que siguen en pie. Un eje de esta campaña ya
produjo exactamente eso, un haz de tres, y la instrumentación actual no sabe
transportarlo porque solo puede leer un umbral de un solo nombre.}

\subsection{La unidad de reporte}
\label{sec:unidad}

\term[def:category]{Category}{La partición por conocimiento previo de la query,
en el momento de la cota inferior de la ventana. Son tres:
\cod{NK} (\emph{no knowledge}), proteínas sin ninguna anotación previa en el
aspecto; \cod{LK} (\emph{limited knowledge}), con anotación en otro aspecto pero
no en este; y \cod{PK} (\emph{prior knowledge}), con anotación previa en este
mismo aspecto. Es la partición que usa CAFA y separa tres problemas distintos:
predecir desde cero, transferir entre aspectos y completar.}

\term[def:aspect]{Aspect}{La sub-ontología: \cod{BPO} (proceso biológico),
\cod{MFO} (función molecular) y \cod{CCO} (componente celular). Son tres grafos
separados con estadística distinta, y promediar entre ellos mezcla tres
problemas.}

\term[def:panel]{Panel}{El producto de category por aspect: nueve celdas.
\textbf{Es la partición obligatoria y nunca se agrega.} Ni la campaña, ni la
superficie que la publica, ni la operación que compara dos arms suman o promedian
los nueve paneles. La operación de comparación se niega explícitamente y nombra
su propio párrafo al negarse.

La razón no es de estilo. Un promedio sin pesos entre poblaciones que difieren en
más de un orden de magnitud es en sí mismo un repesado, y va en la dirección
favorable, porque los estratos más pequeños son los más fáciles. Los nueve
paneles se emiten siempre, con resultado o sin él, para que una corrida parcial no
pueda leerse como completa.

La consecuencia para el reporte es que \textbf{un veredicto es un vector de
nueve componentes y nunca un número}. Cada celda lleva su clave de panel, su
población contada sobre la verdad de la propia ventana y nunca inferida, los
niveles puntuados en orden, y la firmeza del veredicto en esa celda.}

\term[def:estrato]{Stratum}{Un punto en siete ejes, y la unidad por la que se
reporta cuando se quiere más resolución que el panel. Los siete son
\cod{category}, \cod{aspect}, \cod{length}, \cod{homology},
\cod{donor\_evidence}, \cod{taxonomy} y \cod{propagation}.

Los siete no son intercambiables, y la división que importa está escrita en el
código: \cod{category}, \cod{aspect} y \cod{length} son propiedades \emph{de la
query} y se pueden leer siempre; los otros cuatro son propiedades \emph{de una
recuperación} y necesitan un donante efectivamente alineado antes de poder
leerse. Pedir un eje de vecindario sin alineamientos calculados descarta
silenciosamente a la mayoría de las queries y reporta a los supervivientes como
si fuesen la población, que es un error que ya ocurrió.}

\term{Length band}{Banda de longitud de secuencia:
\cod{<=512}, \cod{512-1024}, \cod{1024-2048} y \cod{>2048}. Los cortes son los
límites de contexto de los modelos, no cuantiles del corpus. Una proteína más
larga que el contexto se trunca antes de codificarse, de modo que las bandas
separan secuencias que la representación vio enteras de secuencias que vio en
parte, que es una pregunta distinta de si una proteína es larga. Bandas por
cuantiles mezclarían las dos y se moverían cada vez que se moviese el corpus.}

\term[def:homologia]{Homology band}{Lo cerca que está el donante anotado más
próximo, por identidad de secuencia:
\cod{none}, \cod{<=30}, \cod{30-60}, \cod{60-90} y \cod{>90}.
Los cortes marcan donde la transferencia por similitud cambia de carácter: por
debajo del 30\,\% está la zona de penumbra, donde el alineamiento por sí solo
deja de implicar función compartida; por encima del 90\,\% la transferencia es
casi una copia y dice poco sobre si un método generaliza. \cod{none} es una banda
y no un valor ausente: una query sin ningún donante es una población real y
reportable, y plegarla en la banda más baja escondería los fallos de recuperación
dentro de los resultados de modelado.}

\term{Donor evidence}{Si la anotación del donante más cercano se apoya en
experimento: \cod{exp}, \cod{other} o \cod{none}. Se cruza con la banda de
homología porque las dos juntas son lo que un umbral de homología significa de
verdad. Un donante cercano cuya propia anotación fue inferida por cómputo
transfiere una conjetura, y reportarlo junto a un donante cercano con soporte
experimental promediaría dos afirmaciones distintas en un número.}

\term{Taxonomy band}{A qué distancia está el donante más cercano en el árbol de la
vida: \cod{same}, \cod{close}, \cod{intermediate}, \cod{distant},
\cod{root-only} o \cod{none}. Se resuelve sobre donantes de secuencia distinta a
la query, que es la única lectura bajo la cual \emph{misma especie} es un hecho
sobre el vecindario y no un hecho sobre el protocolo.}

\term[def:propagacion]{Propagation band}{Cuánto tuvo que viajar una anotación para
llegar a la query: \cod{none}, \cod{0}, \cod{0-0.05}, \cod{0.05-0.15} y
\cod{>0.15}. No es la banda de homología. Esa pregunta cuán cerca está el donante
más cercano de cualquier tipo; esta pregunta cuánto más lejos hay que ir para
llegar a un donante cuya propia anotación se apoya en experimento, es decir, la
diferencia entre dos distancias. Una proteína puede estar en un vecindario denso
y aun así lejos de cualquier evidencia experimental, y ese hueco es la distancia
que la anotación recorrió.}

\term[def:floorpop]{Population floor}{El tamaño mínimo que una celda debe
tener para que se le permita votar, derivado de la desviación típica emparejada y
del efecto que se quiere poder detectar. Para el encaminamiento vale
$\lceil (2.8016\,\sigma/0.02)^2 \rceil$. Con un $\sigma$ global de 0.13 da 332,
que es la cifra que la campaña ha usado; con el $\sigma$ implícito de cada panel va
de 120 a 1{,}376, y bajo ese criterio tres de los nueve paneles no alcanzan su
propio suelo. Una celda por debajo del suelo no se reporta como un veredicto débil:
no se reporta.}

\subsection{Los estadísticos}

\term[def:fmax]{\cod{fmax\_w}}{La medida centrada en proteína que usa CAFA. Se
barre un umbral $\tau$, se calculan precisión y exhaustividad promediadas sobre
proteínas y ponderadas por information accretion, y se toma el máximo de la
$F_1$ sobre la curva. \textbf{No admite un error típico de la forma
$\sigma/\sqrt{n}$}, porque es un máximo sobre $\tau$ de una combinación no lineal
de dos medias poblacionales. Cualquier intervalo que se le ponga por esa vía es
inventado.}

\term{\cod{f\_micro\_w}}{La variante agrupada, que acumula los conteos sobre todos
los pares antes de calcular la $F_1$. Da más peso a las proteínas con muchas
anotaciones.}

\term[def:bootstrap]{Bootstrap pareado por proteína}{El estadístico que decide en
esta campaña. Para cada proteína se calcula su propia mejor $F_1$ bajo cada uno de
los dos arms que se comparan, se toma la diferencia por proteína, y se remuestrean
las proteínas con reemplazo para obtener un intervalo de confianza de la media de
esa diferencia. Es \emph{pareado} porque las dos configuraciones se evalúan sobre
las mismas proteínas, lo que elimina la varianza entre proteínas, que es la
dominante.

Este estadístico y \cod{fmax\_w} no son dos formas de medir lo mismo. Premian
comportamientos distintos: \cod{fmax\_w} maximiza una curva agregada y recompensa
predecir poco y con confianza, mientras que el pareado da a cada proteína su
propio óptimo y recompensa cubrir a cada proteína. En los dos ejes que la campaña
ya cerró, los dos estadísticos eligieron ganadores distintos, y en un caso
extremos opuestos de la misma rejilla. De ahí la propuesta de que el estadístico
sea un campo del sello (sección~\ref{sec:marco}).}

\term[def:mde]{MDE, efecto mínimo detectable}{La diferencia más pequeña que una
celda puede distinguir de cero con su población y su varianza. Se lee de la propia
distribución bootstrap y es la cifra que decide si una celda tiene voto. Sobre los
nueve paneles de esta campaña la mediana va de 0.0019 en \cod{PK:BPO} a 0.0256 en
\cod{LK:CCO}, un factor de 13.3; y dentro de un mismo panel varía por un factor de
más de veinte según la comparación, que es la razón por la que se publica como
rango y no como constante.}

\term{Cuantización}{Todo lo almacenado está redondeado a $10^{-4}$. Una diferencia
que se lee como cero desde la base de datos es una \emph{cota superior} del
acuerdo, nunca un cero medido. Junto con el hecho de que el MDE de la celda peor
resuelta es 0.0234, esto significa que una regla basada en estimaciones puntuales
puede decidir a partir de diferencias doscientas veces más pequeñas que lo que el
dato resuelve.}
